Section~\ref{sec:regions} reports every region label and every robustness
index under a critical value calibrated at the pair's own ratio of register
cardinality to training row count. This section gives the construction the
article summarises.

\subsection{The inflation factor has a closed form}

Write the basic interval of Section~\ref{sec:boot} as
$[\hat V - L,\, \hat V + U]$. Inflated by a factor $c$ it becomes
$[\hat V - cL,\, \hat V + cU]$, and it covers the truth $\vartheta$ exactly
when
\[
c \;\geq\; \max\!\left\{\frac{\hat V - \vartheta}{L},\;
\frac{\vartheta - \hat V}{U}\right\}.
\]
The smallest $c$ attaining nominal coverage at level $1 - \alpha$ over a set
of replicates is therefore the $(1-\alpha)$ quantile of that ratio across
them, and its Monte Carlo interval comes from the same order statistics. No
search over candidate factors is needed, and the estimate inherits the
replicate count directly: the plane carries \nPlaneCells\ cells at
\nPlaneReps\ replicates each, over $K/n$ from \planeKnMin\ to \planeKnMax,
on two data-generating worlds rather than one (Table~\ref{tab:plane}).

\subsection{A monotone fit, not a parametric one}

The dependence of $c$ on $K/n$ is well summarised by a log-linear fit whose
slope is \calSlope, with a standard error of \calSlopeSe: the widening a pair
needs grows with its register's cardinality relative to its training set, as
the mechanism of Section~\ref{sec:boot} predicts. That fit is quoted because
it is the interpretable summary, but it is \emph{not} what the corpus is
calibrated with. The parametric fit sits below the measured factor at the
smallest ratios, and using it would under-correct exactly the pairs whose
bands are narrowest --- the pairs where an anti-conservative band does the
most damage. The corpus is therefore calibrated with an isotonic regression
through the measured factors, which is monotone by construction and
interpolates the measurements rather than smoothing through them.
Table~\ref{tab:plane} prints the two curves side by side.

\paragraph{Three candidate curves, on the criterion the monotone one was
adopted under} Round twenty-five's widened plane makes $\log n$ a significant
term (Section~\ref{sec:regions}), so the obvious repair is a two-covariate
parametric fit in $\log(K/n)$ and $\log n$. It is not one. On the plane's
fitted cells the weighted residual sum of squares is \wrssRatioOnly\ for the
ratio-only fit, \wrssWithN\ with $\log n$ added, and \wrssIsotonic\ for the
monotone curve; and on the criterion that matters --- how many cells the
curve sits below their own measured factor's lower confidence limit, which is
the anti-conservative direction --- the two parametric fits are alike at
\nUnderRatioOnly\ and \nUnderWithN\ cells against the monotone curve's
\nUnderIsotonic. Adding $n$ to a parametric curve improves the fit and does
not improve the object. The monotone curve is retained and its residual in
$n$ is reported as a limitation rather than fitted away.

\subsection{What the calibration does and does not correct}

The factor is estimated on a \emph{pointwise} interval in a simulation and
applied to a \emph{simultaneous} band on real logs. It corrects the scale of
the studentised statistic, which is the channel the plane identifies. It does
not correct a change in the shape of the maximum's distribution, and it
assumes that the simulated worlds' dependence on $K/n$ transfers to logs whose
dependence structure and signal are not the simulation's. Because that
transfer is an assumption, Section~\ref{sec:regions} reports the corpus counts
at three settings --- nominal, calibrated, and every factor pushed to the
upper end of its own Monte Carlo interval --- and Section~\ref{sec:lim}
carries the transfer as a limitation.
